% Generated from locale/tr/content/set-theory/choice/tarskiscott.tex; do not hand-edit.


\olfileid[tr]{sth}{choice}{tarskiscott}
\olsection{Tarski--Scott Hilesi}

\olref[cardinals][cardsasords]{defcardinalasordinal} içinde kardinal
sayıları sıra sayıları olarak tanımladık. Bunu yapmak için İyi Sıralama
Aksiyomu'nu varsaydık. Bunu yalnızca ``sektör standardı'' olduğu için
yaptık.

Öyleyse İyi Sıralama'yı çevreleyen felsefi sorunları tartışmadan önce,
sektör standardından \emph{ayrılabileceğimizin} ve İyi Sıralama'yı
\emph{varsaymadan} da bir kardinal sayılar teorisi geliştirebileceğimizin
açık olması önemlidir.
\olref[sfr][siz][]{chap} içinde göründükleri biçimiyle
$\cardeq{A}{B}$, $\cardle{A}{B}$ ve $\cardless{A}{B}$ tanımlarını yine
kullanabiliriz. Yalnızca yeni bir \emph{kardinal sayı} kavramına
ihtiyacımız olacaktır.

Naif bir düşünce, $A$'nın kardinalitesini şöyle tanımlamaya çalışmak olurdu:
\[
	\Setabs{x}{\cardeq{A}{x}}.
\]
Bunu, \olref[cardinals][hp]{sec} kesiminin en sonunda taslağı verilen
Frege'nin $\# x Fx$ tanımıyla karşılaştırabilirsiniz. Orada değindiğimiz
nedenlerle bu tanım başarısız olur. Her tek elemanlı küme
$\{\emptyset\}$ ile eş güçlüdür. Fakat hiyerarşinin her ardıl aşamasında
yeni tek elemanlı kümeler oluşur (yalnızca önceki aşamanın tek elemanlı
kümesini düşünün). Dolayısıyla
$\Setabs{x}{\cardeq{A}{x}}$ var değildir; çünkü bir rankı olamaz.

Bu sorunu aşmak için Tarski ve Scott'a ait bir hile kullanırız:\footnote{Bir
hatırlatma: açıkça tersi belirtilmedikçe bütün formüllerin parametreleri
olabilir.}

\begin{defn}[Tarski--Scott]
Her $\phi(x)$ formülü için $[x:\phi(x)]$, $\phi(x)$ koşulunu sağlayan ve
olanaklı en küçük ranka sahip bütün $x$'lerin kümesi olsun (hiç $\phi$
yoksa $\emptyset$ olsun).
\end{defn}

Bu tanımın meşru olduğunu denetlemeliyiz. $\ZF$ içinde çalışırken
\olref[spine][foundation]{zfentailsregularity}, her $x$ için
$\setrank{x}$'in var olduğunu güvence altına alır. Şimdi $\phi$'yi sağlayan
varlıklar varsa $\alpha$, $(\exists x\subseteq V_\alpha)\phi(x)$ olacak en
küçük rank olsun; yani
$(\forall\beta\in\alpha)(\forall x\subseteq V_\beta)\lnot\phi(x)$ olsun.
O zaman Ayırma ile $[x:\phi(x)]$ kümesini
$\Setabs{x\in V_{\alpha+1}}{\phi(x)}$ olarak tanımlayabiliriz.

Tarski--Scott hilesini gerekçelendirdiğimize göre artık bir kardinalite
kavramını tanımlamak için kullanabiliriz:

\begin{defn}
$A$'nın \textsc{ts}-kardinalitesi
$\text{tsc}(A)=[x:\cardeq{A}{x}]$'dir.
\end{defn}

Bir \textsc{ts}-kardinal sayısının tanımı İyi Sıralama'yı kullanmaz.
Fakat bu Aksiyom olmadan bile \emph{\textsc{ts}-kardinal sayılarının},
\olref[cardinals][cardsasords]{defcardinalasordinal} içinde tanımlanan
\emph{kardinal sayılar} gibi davrandığını gösterebiliriz. Örneğin
\olref[cardinals][cardsasords]{lem:CardinalsBehaveRight} ve
\olref[card-arithmetic][opps]{lem:SizePowerset2Exp} sonuçlarını
\textsc{ts}-kardinal sayıları cinsinden yeniden ifade edersek, ispatlar
İyi Sıralama varsayılmadan $\ZF$ içinde sorunsuzca yürür.

Bu konudayken Tarski--Scott hilesini kullanarak bir sıra sayıları teorisi
de geliştirebileceğimizi belirtmek yerinde olur. $\tuple{A,<}$ bir iyi
sıralama olduğunda
$\text{tso}(A,<)=[\tuple{X,R}:\ordeq{\tuple{A,<}}{\tuple{X,R}}]$ olsun.
Kardinal ve sıra sayılarına bu yaklaşım hakkında daha fazlası için
\citet[chs.~9--12]{Potter2004} eserine bakın.

